\section{Bölüm sonu çözümlü problemler}

\begin{enumerate}
  \item \textbf{Standart puan.} Kaygı puanlarının ortalaması 50, standart sapması 8'dir. 66 puan alan öğrencinin $z$ puanını ve konumunu yorumlayın.

  \textbf{Çözüm:} $z=(66-50)/8=2$'dir. Öğrencinin puanı ortalamanın iki standart sapma üzerindedir. Bu göreli konum tek başına klinik tanı değildir.

  \item \textbf{Normal olasılık.} $X\sim N(100,15^2)$ ise 115'in altında kalma olasılığını bulun.

  \textbf{Çözüm:} $z=(115-100)/15=1$ olur. Standart normal dağılımdan $P(Z<1)\approx0{,}8413$; dolayısıyla olasılık yaklaşık yüzde 84,1'dir.

  \item \textbf{Merkezî limit teoremi.} Evren standart sapması 12 olan bir puan için $n=36$ kişilik örneklem ortalamasının standart hatasını bulun. Evren hafif çarpıksa normal yaklaşımın neden yine kullanılabileceğini açıklayın.

  \textbf{Çözüm:} $SE(\bar X)=12/\sqrt{36}=2$'dir. Bağımsız gözlemler ve sonlu varyans altında örneklem ortalamasının dağılımı, ham puanlar tam normal olmasa da örneklem büyüdükçe normale yaklaşır. Güçlü aykırı değer veya bağımlılık varsa yalnız $n$ değerine güvenilmez.
\end{enumerate}